\chapter{The model and the Abelian property}\label{ch:toppling}

\section{Sand on a grid, and the toppling rule as a table}\label{sec:model}

Sand grains sit on an $L \times L$ square grid. The state at time $t$ is the height
function
\begin{equation}\label{eq:height}
  p_t(\bm{x}) \in \mathbb{N}_0, \qquad
  \bm{x} = (i,j) \in \bm{L} = \{1,\dots,L\}^2 ,
\end{equation}
and the list of all $L^2$ heights is the \emph{configuration}. Time does not advance
during an avalanche: a drop is instantaneous and the whole rearrangement it triggers
belongs to the same time step \cite{handbook2020}. A configuration is \emph{stable}
when every site is below the threshold $k$, and \emph{unstable} otherwise.

An unstable site gives one grain to each of its four nearest neighbours and loses four
grains itself:
\begin{equation}\label{eq:rule}
  p(\bm{x}) \gets p(\bm{x}) - 4, \qquad
  p(\bm{x} \pm \bm{e}_k) \gets p(\bm{x} \pm \bm{e}_k) + 1, \quad k = 1,2 .
\end{equation}
The rule has the four features required for self-organised criticality. The driving is
slow: one grain per drop, at a rate far below the relaxation rate. The threshold is
local: the test $p(\bm{x}) \ge k$ involves one site only. The redistribution is local:
the grains reach the four neighbours and no other site. And the system dissipates: the
grains that would leave the board are deleted. All four are properties of the rule
rather than of a parameter value, so no quantity has to be tuned for the critical state
to appear.

This changes five sites at once, which is awkward to reason about, so it is better
recorded as how much each site $\bm{y}$ changes when $\bm{x}$ topples once. That table
is the \emph{toppling matrix}
\begin{equation}\label{eq:delta}
  \D(\bm{x},\bm{y}) =
  \begin{cases}
    -4, & \bm{x} = \bm{y}, \\
    \phantom{-}1, & \bm{x},\bm{y} \text{ are nearest neighbours}, \\
    \phantom{-}0, & \text{otherwise},
  \end{cases}
\end{equation}
and the rule is recovered by reading row $\bm{x}$: the diagonal entry is what the
toppling site loses, the four non-zero off-diagonal entries are what its neighbours
gain, and the zeros are the sites that do not notice. Summing row $\bm{x}$,
\begin{equation}\label{eq:rowsum}
  \sum_{\bm{y}} \D(\bm{x},\bm{y})
  = -4 + \#\{\text{neighbours of } \bm{x} \text{ inside the board}\} ,
\end{equation}
so the row sum reports conservation and loss in the same number: $0$ for an interior
site, $-1$ at an edge, $-2$ at a corner.

The row sum is negative exactly at the sites that have a neighbour outside the board,
and nowhere else, so the boundary is the only exit for sand and the only reason the
total stays bounded. Without it, on an infinite lattice or with grains that rebound
instead of leaving, every row sum would be zero: the total sand would grow by one with
every drop, no input--output balance would be possible, and both the total mass and the
front of an avalanche would grow without bound, leaving no stationary state and
therefore no exponents. The boundary does not
have to be far away for this to work --- any finite board will do --- which is why an
$L = 50$ board already shows the critical statistics.

\subsection{The four conditions on \texorpdfstring{$\D$}{Delta}}\label{sec:fourconditions}

The handbook isolates four conditions for the dynamics to be well posed
\cite{handbook2020}; they are its equations (4a)--(4d), numbered here
\begin{subequations}\label{eq:fourconditions}
\begin{align}
  \forall \bm{x} \neq \bm{y}: \quad & \D(\bm{x},\bm{y}) = \D(\bm{y},\bm{x}) \ge 0,
    \label{eq:4a} \\[2pt]
  \forall \bm{x}: \quad & \D(\bm{x},\bm{x}) < 0, \label{eq:4b} \\[2pt]
  \forall \bm{x}: \quad & \sum_{\bm{y}} \D(\bm{x},\bm{y}) \le 0, \label{eq:4c} \\[2pt]
  & \sum_{\bm{x},\bm{y}} \D(\bm{x},\bm{y}) < 0 . \label{eq:4d}
\end{align}
\end{subequations}
Each condition constrains one part of the dynamics. Condition~\eqref{eq:4a} says sand
moves between sites without a preferred direction: the non-negativity is what makes
topplings commute (§\ref{sec:commute}), and the symmetry makes $\D$ a real symmetric
matrix, so that $-\D$ is a grounded graph Laplacian and $\det(-\D)$ counts the spanning
trees of the wired graph by the matrix--tree theorem. Condition~\eqref{eq:4b} says a
toppling removes sand from the site that topples; without a strictly negative diagonal
there would be no toppling at all, only redistribution, and the threshold would not be
defined by the matrix. Condition~\eqref{eq:4c} says a toppling never creates sand: it is
\eqref{eq:rowsum} restated, and it is what bounds the total mass. Condition~\eqref{eq:4d}
says at least one place loses sand, which makes $-\D$ a nonsingular $M$-matrix and
therefore makes $\D$ invertible, $\sigma$ unique by \eqref{eq:firing} and every avalanche
finite (Proposition~\ref{prop:termination}).

All four hold for the matrix of \eqref{eq:delta}, and they can be read off it directly.
The off-diagonal entry is $1 \ge 0$ and the neighbour relation is symmetric, so
\eqref{eq:4a} holds. The diagonal entry is $-4 < 0$, so \eqref{eq:4b} holds, and the
threshold is fixed rather than chosen:
\begin{equation}\label{eq:threshold}
  k = -\D(\bm{x},\bm{x}) = 4 .
\end{equation}
The row sum of \eqref{eq:rowsum} takes the values $0$, $-1$ and $-2$ for interior, edge
and corner sites, so \eqref{eq:4c} holds; and since boundary rows are strictly negative
while the board has a boundary, the total in \eqref{eq:4d} is negative. The threshold is
a consequence of conservation, not a free parameter: if the rule sent a number of grains
different from the number of neighbours, the interior row sum would no longer be zero
and \eqref{eq:4c} would fail. That is the constraint which forces the threshold of the
hexagonal lattice from $4$ to $6$ in Chapter~\ref{ch:hexagonal}.

Note the convention: Dhar writes the rule as $z_{\bm{y}} \to z_{\bm{y}} -
\D(\bm{x},\bm{y})$ \cite{dhar1990,dhar1999}, so his $\D$ is the negative of ours and his
conditions are these three with the sign flipped. The sign has to be checked whenever an
exponent is taken from the literature.

\subsection{The nine-site example}\label{sec:example9}

When $L = 3$ the board has nine sites and $\D$ is a $9 \times 9$ table. Numbering sites
row-major, $(i,j) \mapsto 3i+j$, and filling row by row from \eqref{eq:delta}:
\begin{equation}\label{eq:matrix9}
\D =
\begin{array}{c|ccccccccc}
 & 0 & 1 & 2 & 3 & 4 & 5 & 6 & 7 & 8 \\\hline
 0 & -4 &  1 &  0 &  1 &  0 &  0 &  0 &  0 &  0 \\
 1 &  1 & -4 &  1 &  0 &  1 &  0 &  0 &  0 &  0 \\
 2 &  0 &  1 & -4 &  0 &  0 &  1 &  0 &  0 &  0 \\
 3 &  1 &  0 &  0 & -4 &  1 &  0 &  1 &  0 &  0 \\
 4 &  0 &  1 &  0 &  1 & -4 &  1 &  0 &  1 &  0 \\
 5 &  0 &  0 &  1 &  0 &  1 & -4 &  0 &  0 &  1 \\
 6 &  0 &  0 &  0 &  1 &  0 &  0 & -4 &  1 &  0 \\
 7 &  0 &  0 &  0 &  0 &  1 &  0 &  1 & -4 &  1 \\
 8 &  0 &  0 &  0 &  0 &  0 &  1 &  0 &  1 & -4
\end{array}
\end{equation}
with the site indices written outside. Row $4$, the centre, is $(-4,+1,+1,+1,+1)$: the
site takes $-4$ in column $4$, its neighbours $+1$ in columns $1,3,5,7$, the corners
$0$, and the sum is zero as \eqref{eq:rowsum} requires. Row $0$ is a corner, with sum
$-2$, so two grains leave the board whenever it topples.

\subsection{Program: one local function}

Equation~\eqref{eq:rule} is a local function of constant cost whose return value answers
two questions at once: whether the toppling happened, and how many grains left the
board.

\begin{lstlisting}[caption={One toppling. Every line is one entry of row $\bm{x}$ of
\eqref{eq:delta}.}]
ToppleOutcome topple_once(Site x) {
  const Threshold k = mat_.threshold();            // k = 4 = -Delta(x,x)   (2.5b)
  if (h_[idx(x)] < k) return {false, 0};           // stable: T_x is the identity
  h_[idx(x)] -= k;                                 // the diagonal entry
  LossCount loss = 0;
  for (const auto& d : kFourDirections) {
    const auto ni = (long long)x.i + d[0], nj = (long long)x.j + d[1];
    if (ni < 0 || nj < 0 || ni >= n_ || nj >= m_) ++loss;   // row sum < 0 here
    else h_[idx({(std::size_t)ni, (std::size_t)nj})] += 1;  // the +1 entries
  }
  return {true, loss};
}
\end{lstlisting}

The correspondence is item by item: \texttt{mat\_.threshold()} is
$-\D(\bm{x},\bm{x})$, the early return is the branch in which $T_{\bm{x}}$ is the
identity, the subtraction is the diagonal entry, and the loop visits exactly the four
non-zero off-diagonal entries, incrementing \texttt{loss} whenever a neighbour falls
outside the board. So \texttt{loss} is exactly $-\sum_{\bm{y}} \D(\bm{x},\bm{y})$: the
same integer as \eqref{eq:rowsum}, computed the other way round. The two are computed
independently and compared, so that a mistake in the boundary test --- an unsigned index
that wraps around instead of reporting ``outside'' --- cannot corrupt both silently.

\section{Order of toppling, and the stable state}\label{sec:abelianchapter}

An avalanche is the chain of topplings a drop sets off, and it has an order, because the
dynamics says which sites are unstable but never which one to topple next. The rest of
this chapter shows that the choice does not matter and isolates which hypothesis does
the work.

\subsection{The toppling operator, and commutativity}\label{sec:commute}

Write ``topple once at $\bm{x}$'' as an operator:
\begin{equation}\label{eq:Tx}
  (\Tx p)(\bm{y}) =
  \begin{cases}
    p(\bm{y}) + \D(\bm{x},\bm{y}), & p(\bm{x}) \ge k , \\[2pt]
    p(\bm{y}), & p(\bm{x}) < k .
  \end{cases}
\end{equation}
It touches $\bm{x}$ and its four neighbours and acts only when the site should topple;
otherwise it is the identity. Toppling $\bm{x}$ $n$ times adds row $\bm{x}$ $n$ times,
$\Tx^{\,n} p = p + n\,\D(\bm{x},\cdot)$, and a whole avalanche is a product
$\stab = \prod_i T_{\bm{x}_i}$ in the order the topplings happened.

\begin{lemma}\label{lem:commute}
  Let $\bm{x}_1 \neq \bm{x}_2$ satisfy $p(\bm{x}_1) \ge k$ and $p(\bm{x}_2) \ge k$.
  Then both $T_{\bm{x}_2} T_{\bm{x}_1} p$ and $T_{\bm{x}_1} T_{\bm{x}_2} p$ are defined,
  and they are equal.
\end{lemma}

\begin{proof}
  Put $q = T_{\bm{x}_1} p$. By \eqref{eq:4a}, $\D(\bm{x}_1,\bm{x}_2) \ge 0$, so
  $q(\bm{x}_2) = p(\bm{x}_2) + \D(\bm{x}_1,\bm{x}_2) \ge p(\bm{x}_2) \ge k$: toppling
  $\bm{x}_1$ cannot calm $\bm{x}_2$ down, because a toppling only adds sand to other
  sites. Hence the branch $q(\bm{x}_2) \ge k$ applies and
  $\bigl(T_{\bm{x}_2} T_{\bm{x}_1} p\bigr)(\bm{y}) = p(\bm{y}) + \D(\bm{x}_1,\bm{y})
  + \D(\bm{x}_2,\bm{y})$. The same computation with the roles exchanged has the same
  right-hand side, because addition of integers in $\mathbb{Z}^{\bm{L}}$ is
  commutative.
\end{proof}

The proof uses $\D(\bm{x}_1,\bm{x}_2) \ge 0$ and nothing else: not the symmetry half of
\eqref{eq:4a}, not the row sums, not the number of neighbours, not the shape of the site
set. That is why the result survives a change of lattice. It is not the only hypothesis
at work, however.

\begin{remark}\label{rem:withdrawal}
  Equation~\eqref{eq:Tx} asks two things of a rule: that the instability test at
  $\bm{x}$ involves the height of $\bm{x}$ alone, and that a toppling only adds sand to
  other sites. The handbook's gradient rule fails the first, since $\bm{x}$ is unstable
  when it exceeds a neighbour by $k$ and the trigger therefore involves $p(\bm{y})$; a
  rule that may \emph{take} a grain from a neighbour fails the second. Neither property
  is equivalent to determinism, and random rules are not excluded: the Manna model
  \cite{manna1990}, which hands grains to randomly chosen neighbours, is still Abelian
  \cite{dhar1999}, because the randomness is attached to each site in advance. Randomness
  drawn \emph{at the moment of toppling} is different, and Dhar's non-Abelian variant
  \cite{dhar1999} loses the property that way.
\end{remark}

\subsection{Termination and uniqueness}\label{sec:termination}

\begin{proposition}\label{prop:termination}
  Let $\D$ satisfy \eqref{eq:4a}--\eqref{eq:4d} on the finite board $\bm{L}$. Then every
  legal sequence of topplings from $p$ has at most $\bm{1}^{\!\top}(-\D)^{-1}p$ terms.
\end{proposition}

\begin{proof}
  Let $v_t \ge 0$ count the topplings in the first $t$ steps; then $p_t = p +
  \D^{\!\top} v_t$, and since height is non-negative, $(-\D)v_t \le p$ using the
  symmetry $\D^{\!\top} = \D$. Now $-\D$ has positive diagonal by \eqref{eq:4b},
  non-positive off-diagonal by \eqref{eq:4a} and non-negative row sums by \eqref{eq:4c},
  strictly positive for at least one row by \eqref{eq:4d}, with the off-diagonal pattern
  of the connected grid graph. It is therefore irreducibly diagonally dominant, hence a
  nonsingular $M$-matrix \cite{berman1994}, so $(-\D)^{-1} \ge 0$ entrywise and
  $0 \le v_t \le (-\D)^{-1}p$ for every $t$. Each toppling increases
  $\sum_{\bm{x}} v_t(\bm{x})$ by one, so the number of steps is bounded by the stated
  constant, which depends only on the starting configuration.
\end{proof}

\begin{theorem}[Abelian lemma]\label{thm:abelianlemma}
  Let $v$ be the count vector of a legal sequence taking $p$ to a stable configuration,
  and let $w \ge 0$ be any integer vector with $p + \D^{\!\top} w$ stable. Then
  $v \le w$ componentwise.
\end{theorem}

\begin{proof}
  Suppose not, and let $t$ be the first step at which some site topples for the
  $(w(\bm{x})+1)$-st time. Before step $t$ every site has toppled at most $w$ times, so
  $v_{t-1} \le w$ with equality at $\bm{x}$. Step $t$ is legal, so
  $p_{t-1}(\bm{x}) \ge k$; but
  \begin{equation}\label{eq:contradiction}
    p_{t-1}(\bm{x})
      = p(\bm{x}) + w(\bm{x})\D(\bm{x},\bm{x})
        + \sum_{\bm{y} \neq \bm{x}} v_{t-1}(\bm{y})\, \D(\bm{y},\bm{x})
      \le p(\bm{x}) + \bigl(\D^{\!\top} w\bigr)(\bm{x}) \ \le\ k-1 ,
  \end{equation}
  since $\D(\bm{y},\bm{x}) \ge 0$ for $\bm{y} \neq \bm{x}$ lets $v_{t-1} \le w$ be
  pushed through the sum term by term, and $p + \D^{\!\top}w$ is stable. That
  contradiction proves $v \le w$.
\end{proof}

\begin{theorem}\label{thm:unique}
  Every complete legal sequence from $p$ produces the same final configuration and the
  same toppling-count vector.
\end{theorem}

\begin{proof}
  Every sequence is finite by Proposition~\ref{prop:termination}. For two complete
  sequences with count vectors $v, v'$, Theorem~\ref{thm:abelianlemma} applied to $v$
  with $w = v'$ gives $v \le v'$, and applied to $v'$ with $w = v$ gives $v' \le v$.
  Hence $v = v'$ and the endpoints, being $p + \D^{\!\top}v$, coincide.
\end{proof}

So the stabilisation map $p \mapsto p^\circ$ is well defined, an avalanche has a
definite size, and the simulation may topple in whatever order is convenient. The count
vector is determined by the endpoints alone,
\begin{equation}\label{eq:firing}
  p^\circ = p + \D^{\!\top}\sigma
  \qquad\Longleftrightarrow\qquad
  \sigma = \bigl(\D^{-1}\bigr)^{\!\top} \bigl(p^\circ - p\bigr),
\end{equation}
with no sink vertex or row deletion needed, since \eqref{eq:4d} already makes $\D$
invertible. This is the \emph{least action principle} \cite{fey2010,klivans2019}:
$\sigma$ is the pointwise smallest non-negative integer vector for which
$p + \D^{\!\top}\sigma$ is stable.

\subsection{Four observables, and the program}\label{sec:observables}

Once $\sigma$ is known the avalanche is completely described, and the four quantities
the handbook asks for \cite{handbook2020} are four summaries of it:
\begin{align}
  \text{topples } s &= \sum_{\bm{x}} \sigma(\bm{x}), &
  \text{area } s_d &= \abs{\{\bm{x} : \sigma(\bm{x}) > 0\}}, \notag \\[2pt]
  \text{loss } l &= \sum_{\bm{x}} \sigma(\bm{x}) \Bigl(-\sum_{\bm{y}}\D(\bm{x},\bm{y})\Bigr), &
  \text{length } R &= \max_{\bm{y}:\,\sigma(\bm{y})>0} d(\bm{x}_0, \bm{y}) .
  \label{eq:observables}
\end{align}
The first two count topplings with and without repetition, so $s_d \le s$, with equality
exactly when no site topples twice. The loss can only be positive if the avalanche
reached the boundary, since an interior toppling has row sum zero. The length is the
graph distance from the drop site to the farthest site that toppled; because
\eqref{eq:rule} only moves grains along the four grid directions, that distance is the
Manhattan distance
\begin{equation}\label{eq:manhattan}
  R = \max_{\bm{y}:\,\sigma(\bm{y})>0}
      \bigl(\abs{i_{\bm{y}} - i_0} + \abs{j_{\bm{y}} - j_0}\bigr).
\end{equation}
The choice of metric belongs to the geometry, not to the reporting: $R$ counts steps, so
the metric must be the one the grains actually walk, and on the hexagonal lattice of
Chapter~\ref{ch:hexagonal} it is not the Manhattan distance.

``Topple until stable'' is a queue loop: unstable sites enter, are toppled one at a
time, and any site made unstable by a toppling enters in turn.

\begin{lstlisting}[caption={One drop and the relaxation. The order inside the queue is
arbitrary by Theorem~\ref{thm:unique}.}]
void drop_sand_at(Site s, std::uint32_t n = 1) {
  h_[idx(s)] += n;
  if (h_[idx(s)] >= mat_.threshold()) stabilize(s);   // avalanche
  else history_.push_back(AvalancheStats{});          // no avalanche: zeros
}

AvalancheStats stabilize(Site start) {
  std::deque<Site> q{start};
  std::vector<char> in_q(n_ * m_, 0);                 // dedup on entry
  in_q[idx(start)] = 1;
  ToppleCount s = 0; LossCount loss = 0; std::size_t area = 0, R = 0;
  while (!q.empty()) {
    const Site x = q.front(); q.pop_front(); in_q[idx(x)] = 0;
    if (h_[idx(x)] < k) continue;                     // already absorbed
    loss += topple_once(x).loss; ++s;
    for (const auto& d : kFourDirections) { /* push unstable neighbours */ }
  }
  for (const Site y : fired_list_) R = std::max(R, dist_l1(y, start));
  return {s, area, loss, R};
}
\end{lstlisting}

Three details of the loop bear on the theorems. The height of a site is tested again
when it leaves the queue, because a neighbour's toppling may have added sand to it while
it waited; this is simpler than maintaining the invariant that every queued site is
unstable, and it is safe because by \eqref{eq:4a} an unstable site stays unstable until
it is toppled. Deduplication on entry changes only the queue length. And replacing the
queue by a stack or by a random order changes nothing about the result: that is the
Abelian property, and the two-order comparison is the cheapest test of it.
